\documentclass[11pt,letterpaper]{article}
\usepackage[margin=1in]{geometry}
\usepackage{amsmath,amssymb,amsthm}
\usepackage{graphicx}
\usepackage{hyperref}
\usepackage{booktabs}
\usepackage{enumitem}
\usepackage{bm}
\usepackage{xcolor}

\hypersetup{
  colorlinks=true, linkcolor=blue!60!black, citecolor=blue!60!black,
  urlcolor=blue!60!black,
}

\providecommand{\Systrophe}{Systroph\=e}

\title{\Large\textbf{Distance-5 rotated surface code break-even crossing on \texttt{ibm\_kingston}}\\
\large{MWPM-decoded $d=5$ Z-memory at $n_{\mathrm{rounds}} = 1$ exceeds the wall-clock-matched bare-qubit baseline by $+2.88$\,pp $= 4.0\sigma$}}

\author{Christian Knopp \\
  \texttt{cknopp@gmail.com} \\
  \href{https://github.com/Zynerji/systrophe}{\texttt{github.com/Zynerji/systrophe}}}

\date{2026-05-12 \quad / \quad \Systrophe{}\ v0.19.x post-tag}

\begin{document}
\maketitle

\begin{abstract}
We report a \textbf{sustained} logical-versus-physical break-even crossing for the distance-5 rotated surface code's Z-memory on IBM Quantum's \texttt{ibm\_kingston} 156-qubit Heron-r2 processor. The experiment encodes $|0_L\rangle$ on a $5 \times 5$ data-qubit grid (25 data + 12 Z-stabilizer ancillas = 37 qubits), performs $n_{\mathrm{rounds}} \in \{1, 2, 4\}$ rounds of Z-syndrome extraction, and decodes the syndrome difference history via minimum-weight perfect matching (MWPM) with full Dijkstra-shortest-path X-flip corrections on the dual-lattice graph. The Dijkstra-MWPM-decoded logical recovery rate beats the wall-clock-matched bare-qubit baseline at every round count tested:
\begin{itemize}[leftmargin=*,itemsep=0pt]
\item $n=1$: logical $0.9316$ vs bare $0.8528$, diff $+0.079$ ($\bm{11.6\sigma}$)
\item $n=2$: logical $0.8167$ vs bare $0.7754$, diff $+0.041$ ($4.6\sigma$)
\item $n=4$: logical $0.7869$ vs bare $0.7024$, diff $+0.085$ ($\bm{8.8\sigma}$)
\end{itemize}
This is the first \textbf{multi-round} logical-error-rate below physical-error-rate demonstration on Heron-r2 in our \Systrophe{} program, and supersedes the $d=3$ Steane sub-threshold baseline reported in \texttt{paper/steane\_logical\_qubit.pdf}. The earlier single-qubit-flip approximation of MWPM yielded only a single-round crossing; switching to a proper Dijkstra-shortest-path correction chain on the dual-lattice graph recovered $5$--$10$ percentage points of additional logical fidelity at every round count and sustained the crossing. We still do \emph{not} claim to match the SOTA QEC threshold demonstrations of Google Quantum AI Willow (distance-7, fault-tolerant) \cite{willow} or IBM bivariate-bicycle codes \cite{bb_paper}; we report a clean and reproducible sustained-break-even crossing on Heron-r2 with a documented MWPM pipeline.
\end{abstract}

\section{Introduction}

The defining hardware demonstration in quantum error correction is the \emph{break-even} crossing: at matched wall-clock time, the encoded logical qubit retains $Z$ (or $X$) information better than a bare physical qubit. Below this crossing the QEC overhead is net-harmful; above it, encoding pays off. The classical milestone is Google Quantum AI's distance-7 Willow result (Dec 2024) \cite{willow} showing sustained below-threshold scaling. Earlier IBM and Google distance-3 to distance-5 comparisons found a fragile crossing only at distance-5, with very specific decoder choices and round counts \cite{willow_precursor}.

In our companion paper \cite{steane_paper} we reported a $d = 3$ Steane \texttt{[[7,1,3]]} round-sweep experiment on \texttt{ibm\_kingston} with $T_{1,L} / T_{1,\mathrm{phys}} = 0.579$ --- a clean sub-threshold result. This paper reports the natural follow-up: a $d = 5$ rotated surface code on the same chip with the same wall-clock-matched baseline methodology. We find a single-round break-even crossing at $4.0\sigma$, but no crossing at higher round counts within our naive MWPM implementation.

\section{Experiment}

\subsection{Code and circuit}

We use the canonical rotated $d = 5$ surface code with rough top/bottom boundaries (X-edges) and smooth left/right boundaries (Z-edges) on a $5 \times 5$ data-qubit grid (25 data qubits). The Z-stabilizers comprise eight weight-4 plaquettes plus four weight-2 boundary half-plaquettes, totalling $N_Z = 12$. We additionally allocate 12 Z-ancilla qubits for syndrome extraction (one per Z-stab), giving a total qubit count of 37.

For the Z-memory experiment, encoding $|0_L\rangle$ is trivial: the data register initialised to $|0\rangle^{25}$ is automatically the $+1$ eigenstate of all Z-stabilizers and has $Z_L = +1$. No explicit encoder circuit is required when only Z-syndromes are extracted (which suffices for protection against X-type errors). Each round of Z-syndrome extraction applies, for each Z-stab, CNOTs from each data qubit in the stab's support into the dedicated ancilla, followed by ancilla measurement and reset.

We do not extract X-syndromes in this experiment, so X-type errors propagating into Z observables during measurement are not corrected. This is a deliberate trade-off: the Z-memory experiment is the simplest valid $d = 5$ surface-code memory test, and the X-syndrome extension is the natural next step.

\subsection{Submission}

The full round sweep was submitted as a single SamplerV2 batch (job \texttt{d81kel6gbeec73akmp00}, 4096 shots/circuit) on \texttt{ibm\_kingston} with:
\begin{itemize}[leftmargin=*,itemsep=0pt]
\item Optimization level 3, ISA-transpiled to Heron-r2 heavy-hex topology.
\item Dynamical decoupling enabled (XpXm sequence).
\item Three Z-memory circuits at $n_{\mathrm{rounds}} \in \{1, 2, 4\}$, paired with three bare-qubit baselines at matched compiled durations $\Delta t$ (using \texttt{qc.delay(dt)} on a single qubit + measurement).
\end{itemize}

ISA depths and durations:
\begin{center}
\begin{tabular}{cccc}
\toprule
$n_{\mathrm{rounds}}$ & Surface ISA depth & Duration ($\mathrm{dt}$ units) & Bare ISA depth \\
\midrule
1 & 152 & 30\,400 & 2 \\
2 & 289 & 57\,800 & 2 \\
4 & 465 & 93\,000 & 2 \\
\bottomrule
\end{tabular}
\end{center}

Note that the $d = 5$ surface code's ISA depth at $n_{\mathrm{rounds}} = 1$ (152) is \emph{shallower} than our $d = 3$ Steane code at the same round count (229), because the surface code's Z-syndrome circuit is parallel-friendly (no $H$-wrapping required) and the heavy-hex transpiler routes the surface code's local nearest-neighbour CNOTs more efficiently than the Steane code's non-local stabilizer support.

\subsection{Minimum-weight perfect matching decoder}

The decoder takes the syndrome difference history across rounds plus the final data-derived Z-syndrome and decodes via \texttt{networkx.min\_weight\_matching} on a graph with one node per violated stabilizer per round. Edges are weighted by lattice-Manhattan distance between stab centers. For each matched pair, we apply $X$-flip corrections to the shortest-path data qubit (or the shared-support qubit if any). Unpaired (odd-cardinality) violations are mapped to their first-support qubit.

This is an approximate MWPM: a fully fault-tolerant implementation would use a 3D syndrome graph (round × stab × stab) with Dijkstra-shortest-path edge weights and explicit boundary-edge nodes. Our simplified implementation suffices to demonstrate the single-round break-even crossing reported here.

We compare the MWPM decoder to a baseline ``greedy'' decoder that uses only the final data-derived Z-syndrome and applies a fixed correction per violated stab. The greedy decoder is essentially the same as the $d = 3$ Steane lookup-table decoder; for $d = 5$ it cannot resolve multi-qubit error chains and performs worse than MWPM.

\section{Results}

\begin{center}
\begin{tabular}{rrrrrrr}
\toprule
$n_{\mathrm{rounds}}$ & Greedy & Approx-MWPM & \textbf{Dijkstra-MWPM} & Bare & D-MWPM $-$ bare & Significance \\
\midrule
1 & $0.6990$ & $0.8816$ & $\bm{0.9316}$ & $0.8528$ & $\bm{+0.0789}$ & $\bm{11.6\sigma}$ \\
2 & $0.6013$ & $0.7429$ & $\bm{0.8167}$ & $0.7754$ & $\bm{+0.0413}$ & $4.6\sigma$ \\
4 & $0.5337$ & $0.6821$ & $\bm{0.7869}$ & $0.7024$ & $\bm{+0.0845}$ & $\bm{8.8\sigma}$ \\
\bottomrule
\end{tabular}
\end{center}

The headline numbers:
\begin{itemize}[leftmargin=*,itemsep=0pt]
\item \textbf{The Dijkstra-MWPM-decoded $d = 5$ surface code Z-memory beats the wall-clock-matched bare-qubit baseline at every round count tested, at $4.6\sigma$ to $11.6\sigma$ significance.}
\item The greedy ``flip first qubit'' decoder gives only $0.70$ at $n=1$ and $0.53$ at $n=4$ --- well below bare. Decoder choice is the dominant factor.
\item Approximate MWPM (single-qubit-flip per matched pair, the first version) gives a single-round crossing of $+2.88$ pp ($4.0\sigma$) but loses at $n \ge 2$. The full Dijkstra-shortest-path version (flipping ALL data qubits along the matched lattice path) recovers $5$--$10$ pp at every round count and sustains the crossing.
\item At $n = 4$ the gap is $+8.45$ pp at $8.8\sigma$ --- larger than at $n = 1$ in percentage-point terms, because the bare-baseline decays faster than the logical does. This is the canonical QEC signature: encoded logical lifetimes scale with distance and decoder quality, bare lifetimes do not.
\end{itemize}

\subsection{Why the full Dijkstra-MWPM sustains the crossing}

The single-qubit-flip MWPM (our first version) used to fail past $n = 1$ because longer error chains were under-corrected: a matched pair separated by lattice distance $k$ in the dual lattice would be corrected by flipping only ONE data qubit, leaving the other $k - 1$ qubits in the chain still flipped. With increasing round count and accumulating errors, this under-correction compounded.

The Dijkstra-MWPM version finds the proper shortest path through the stab-adjacency graph (where adjacent Z-stabs share a data qubit) for each matched pair, then flips ALL data qubits along the path. This is the canonical surface-code correction recipe, and it recovers the full $d = 5$ protection at every round count we tested.

We still expect saturation at large $n$ because (a) the decoder uses syndrome difference-history per round, not a full 3D syndrome graph; (b) Z-syndrome-only memory does not detect phase errors on the ancillas; (c) accumulated correlated errors eventually exceed the code's distance-5 correction radius. Pushing past these limits is the natural next step ($d = 7$, full X+Z syndromes, 3D MWPM).

\subsection{Comparison to literature SOTA}

Google Quantum AI's distance-3 vs distance-5 surface code comparison in 2023 \cite{willow_precursor} reported a similar single-round-style crossing that disappeared at higher rounds without further architectural improvements. The Willow result \cite{willow} achieved sustained sub-threshold scaling at distance-7 with much higher gate fidelities (cz error $\sim 5 \times 10^{-4}$ vs our Heron-r2 $\sim 2 \times 10^{-3}$). Our result is consistent with the literature picture: $d = 5$ on current Heron-r2 gate fidelities lives near the threshold; with proper MWPM at single-round measurement we cross break-even; sustained crossing across rounds requires either gate-fidelity improvement (Willow-class) or larger distance ($d \ge 7$).

\textbf{Honest claim-level (Systroph\=e taxonomy)}:
\begin{itemize}[leftmargin=*,itemsep=0pt]
\item Reproducible: yes; rerun \texttt{python experiments/surface\_code\_d5\_mwpm.py} or resubmit via \texttt{surface\_code\_d5\_round\_sweep.py --submit --n-rounds 1}.
\item Novel demonstration on Heron-r2: yes; first single-round $d = 5$ surface-code break-even crossing in our experimental program.
\item SOTA contribution: no; this does not improve on Willow or BB-code published thresholds.
\end{itemize}

\section{Open extensions}

\begin{enumerate}[leftmargin=*,itemsep=0pt]
\item \textbf{Full fault-tolerant MWPM} with 3D syndrome graphs and boundary nodes (Dennis-Kitaev-Landahl-Preskill 2002).
\item \textbf{X + Z syndrome extraction} (full surface code memory, not just Z).
\item \textbf{Higher distance} ($d = 7$ would need 49 data + 24 ancillas = $\sim 90$ qubits, still fits Heron-r2's 156).
\item \textbf{More shots} at $n_{\mathrm{rounds}} = 1$ to push the break-even-crossing significance from $4.0\sigma$ to $10\sigma+$.
\item \textbf{Multi-logical-qubit demonstration}: parallel patches of the same code on disjoint Heron-r2 subgraphs.
\end{enumerate}

\section{Reproducibility}

All code and data are at \href{https://github.com/Zynerji/systrophe}{\texttt{github.com/Zynerji/systrophe}}:

\begin{itemize}[leftmargin=*,itemsep=0pt]
\item Circuit construction: \texttt{experiments/surface\_code\_d5.py}
\item Round sweep submission: \texttt{experiments/surface\_code\_d5\_round\_sweep.py}
\item MWPM decoder: \texttt{experiments/surface\_code\_d5\_mwpm.py}
\item Job ID: \texttt{d81kel6gbeec73akmp00} (round sweep + bare baselines)
\item Raw analysis JSON: \texttt{experiments/results/surface\_code\_d5\_mwpm\_analysis.json}
\end{itemize}

\begin{thebibliography}{9}

\bibitem{willow} Google Quantum AI. \emph{Quantum error correction below the surface code threshold}. Nature \textbf{636}, 456 (Dec 2024).
\bibitem{willow_precursor} Google Quantum AI. \emph{Suppressing quantum errors by scaling a surface code logical qubit}. Nature \textbf{614}, 676 (2023).
\bibitem{bb_paper} S.~Bravyi et al. \emph{High-threshold and low-overhead fault-tolerant quantum memory}. Nature \textbf{627}, 778 (2024).
\bibitem{steane_paper} C.~Knopp. \emph{Distance-3 Steane code repeated-round logical-qubit experiment on \texttt{ibm\_kingston}}. \Systrophe{} repository, \texttt{paper/steane\_logical\_qubit.pdf} (2026).

\end{thebibliography}

\end{document}
